% Independent proof audit; includable in the final report.
\section{Computation sources certified by construction}
\label{sec:constructed-supply-audit}

The computation packet has a mode field $e$ and a triple $[p,a,r]$.
The old packet interface already had an unused field in its computation
branch; that field now carries $e$.  The kernel generates the source
counter itself.  The relevant code is
\begin{verbatim}
    n, t, r = s
    h = [[u and [n], t, r], i(r, s, 1) if u else t]
\end{verbatim}
Here \texttt{u} is the mode $e$, \texttt{s} is $[p,a,r]$, and
\texttt{t} is reassigned to the initializer $a$.  Canonical mode values
are $0$ for the base equation and $1$ for the successor equation.

For base mode the generated equation is
\[
                 I(0,a;v.r)=a.
\]
The submitted predecessor payload $p$ is unused in this case.  For step
mode the generated equation is
\[
                 I(Sp,a;v.r)=r[I(p,a;v.r)/v].
\]
The step counter is a singleton array because the kernel constructs
\texttt{[n]}.  It is never an untrusted purported successor.  Thus no
counter-shape validator is needed in this branch.

The expression \texttt{u and [n]} can return a noncanonical falsy mode
unchanged.  This causes no additional arithmetic case.  If the subsequent
common opening succeeds, a falsy numerical scalar or empty array denotes
zero.  Unsupported falsy JSON values reject when traversed.  A truthy
mode is used only as a branch tag and always produces the same singleton
constructor.  Its own JSON shape has no numerical meaning.

\subsection{Demanded images}

The opener validates an inserted image at numerical root whenever a
meaningful exposed numerical coordinate requests that image.  Consequently,
if opening $r$ demands the predecessor image $[p,a,r]$, the latter is
checked before insertion.  This validation is exactly the old eager
check, performed at the first point where its result can affect the
equation.  Numerical-root validation is idempotent on successful results:
structural recursion leaves permitted coordinates unchanged, rejects a
missing meaningful numerical image, and replaces only inessential
formula leaves in checked annotations.

There is one apparent binding mismatch to check.  In the constructed
source, the displayed successor contributes an additional dummy zero
slot to the environment for $p$.  The predecessor image has only the
iterator's usual two dummy slots.  Root validation of the predecessor
image excludes every meaningful reference to the extra outer slot.
All genuine internal slots and existing rigid names therefore have the
same meanings in both occurrences of $p$.  The update's two aliased
state slots are replaced by the validated predecessor image.  After
decoding, this gives precisely the scalar successor equation.

\subsection{Unused images}

If opening $r$ never requests its image, $r$ is independent of its outer
state on every meaningful numerical path.  Indeed, at opening depth
$d$, references to internal binders are below $d$, while references to
either removed state slot would request the image.  Paths beneath a
numerical binary annotation have no numerical meaning and do not affect
this argument.  Existing rigid names remain parameters.

The common continuation still opens the complete generated equation,
and in particular traverses its source iterator.  It therefore validates
the submitted predecessor, initializer, and update in the environments
where they actually occur.  A malformed unused predecessor image is
not itself a term used by the equation.  Let $w$ be the newly allocated
parameter.  After setting $w=0$ and completing exposed outer slots by
zero, write the source data as $p_0,a_0,r_0$.  The generated equation is
then
\[
 I(Sp_0,a_0;v.r_0)=r_0,
\]
with $r_0$ independent of $v$.  The scalar successor computation axiom
proves it, because substitution of the unused recursive value changes
nothing.  Neither a validation nor a dummy-slot comparison of that
unused raw recursive image is required.

\subsection{The base equation and the fresh parameter}

In base mode the source initializer occurs under two dummy zero slots,
whereas its copy on the right may expose those coordinates during the
common opening.  The standard zero-extension argument applies: substitute
zero for the new parameter, preserve genuine local binders and old rigid
names, and compare the two finite syntax trees structurally.  Their
decoded initializers agree.  The zero instance is therefore the base
computation equation.

If the previous name floor is $k_0$, the new name is $w=k_0-1$.
The original goal and retained assumptions were checked before this
allocation, so they have no meaningful occurrence of $w$.  Generated
source data may acquire $w$ at the later common opening.  The three
cases above establish a computation proof of the actual added equation's
zero instance, $H(0)$.  A successful continuation
\(\Gamma,H(w)\vdash A\) is discharged by renaming its later
eigenparameters, substituting $w=0$, and cutting with that computation
proof.  The old end sequent is unchanged.  The rule never asserts that
an arbitrary added equation $H(w)$ holds for every $w$.

\subsection{Canonical completeness without zero aliases}

Every ordinary base computation is represented by mode $0$, its actual
initializer and update, and an arbitrary unused predecessor payload
(canonically zero).  Every ordinary successor computation with
predecessor $p$ is represented by mode $1$ and the triple $[p,a,r]$.
In particular, a genuine zero predecessor uses \texttt{n=0} together
with \texttt{u=1}, and the generated source counter is exactly
\texttt{[0]}.  No truthy zero alias or additional equality proof is
needed.

As in the inherited canonical emitter, active passive parameters are
allocated rigid names, and dummy slots contain zero rather than
references to otherwise unused enclosing slots.  Thus inserting the
displayed successor introduces no new dependency, and canonical opened
updates have exactly their specified state substitutions.  The result
is the original computation equation in the original raw canonical
syntax.  Supplying it and immediately applying the hypothesis rule
simulates an ordinary computation leaf.  Other proof instructions and
their completeness transformations are unchanged, apart from the
already specified fresh-label accounting.

\paragraph{Precise scope.}
This is a new certificate format and a conservative computation generator.
An old packet whose source counter was malformed is not being newly
accepted with its old meaning.  The new packet supplies a predecessor,
and the checker constructs a genuine successor around it.  Its actual
generated equation is the one that must be used in the continuation.
For example, a binary numerical predecessor denotes zero and yields a
successor counter; it cannot trigger the false equation obtained by
mistaking that binary node itself for a successor.  The entire argument
is about finite syntax and ordinary computation axioms on the full
numerical domain.  It does not use numeral testing as a substitute for
the theorem, and does not claim an absolute smallest proof checker.
